\section{Lösen von Differentialgleichungen}

Diese Section zeigt ein \textbf{systematisches Vorgehen} zum Lösen typischer
Differentialgleichungen in der Elektrotechnik und Leistungselektronik.
Fokus: \textbf{1. Ordnung} (R--L / R--C) und \textbf{2. Ordnung} (RLC),
inkl. \textbf{Stückweise-Betrieb} (Schaltzustände), wie er bei Umrichtern ständig vorkommt.

\smallskip

\textbf{Zielbild:}
\begin{itemize}
    \item DGL-Typ schnell erkennen
    \item Standardlösung aufbauen (homogen + particular)
    \item Anfangs- und Randbedingungen korrekt einsetzen
    \item Laplace-Methode sicher anwenden
    \item Stückweise Lösungen sauber verketten
\end{itemize}

% -------------------------------------------------
\subsection{DGL-Typen (Prüfungs-Screening)}

\subsubsection{1. Ordnung (typisch R--L, R--C)}

\[
\cbl{a \frac{dy}{dt} + b y = g(t)}
\]

\subsubsection{2. Ordnung (typisch RLC, mechanische Systeme)}

\[
\cbl{a \frac{d^2y}{dt^2} + b \frac{dy}{dt} + c y = g(t)}
\]

\subsubsection{Stückweise DGL (Umrichterbetrieb)}

\[
\cbl{\text{Je Schaltzustand eigene DGL, dann Übergangsbedingungen.}}
\]

\crd{\textrightarrow\ 80\% der Punkte kommen aus sauberem Strukturieren, nicht aus Rechenkunst.}

% -------------------------------------------------
\subsection{Universales Vorgehen}

\begin{enumerate}
    \item \textbf{Signal und System definieren:} gesuchte Grösse $y(t)$, Eingänge, Parameter.
    \item \textbf{DGL aufstellen:} KVL/KCL oder Energie-/Kraftgleichungen.
    \item \textbf{DGL normieren:} auf Standardform bringen.
    \item \textbf{Lösungsansatz wählen:}
    \begin{itemize}
        \item Zeitbereich (homogen + particular)
        \item Laplace (algebraisch)
    \end{itemize}
    \item \textbf{Konstanten bestimmen:} Anfangswerte, Periodenbedingungen, Randbedingungen.
    \item \textbf{Plausibilisieren:} Grenzfälle $t\to 0$, $t\to\infty$, Einheiten, Vorzeichen.
\end{enumerate}

% -------------------------------------------------
\subsection{DGL 1. Ordnung im Zeitbereich}

\subsubsection{Standardform}

\[
\cbl{\frac{dy}{dt} + \frac{1}{\tau} y = k \, u(t)}
\qquad
\cbl{\tau = \frac{a}{b}}
\]

\subsubsection{Lösungsstruktur}

\[
\cbl{y(t) = y_h(t) + y_p(t)}
\]

Homogene Lösung ($u(t)=0$):
\[
\cbl{y_h(t) = C \e^{-t/\tau}}
\]

Partikuläre Lösung hängt von $u(t)$ ab.

\subsubsection{Wichtiger Spezialfall: Konstantanregung $u(t)=U$}

Partikulär:
\[
\cbl{y_p = k U \tau \cdot \frac{1}{\tau} = kU}
\]

Gesamtlösung:
\[
\cbl{y(t) = kU + \big(y(0)-kU\big)\e^{-t/\tau}}
\]

\crd{\textrightarrow\ Merksatz: Endwert + (Start-Endwert)$\cdot e^{-t/\tau}$.}

% -------------------------------------------------
\subsection{Beispiel 1: R--L-Kreis (Stromanstieg)}

KVL:
\[
\cbl{L\frac{di}{dt} + Ri = U}
\qquad
\Rightarrow
\qquad
\frac{di}{dt} + \frac{R}{L} i = \frac{U}{L}
\]

Zeitkonstante:
\[
\cbl{\tau = \frac{L}{R}}
\]

Lösung:
\[
\cbl{i(t) = \frac{U}{R} + \left(i(0) - \frac{U}{R}\right)\e^{-t/\tau}}
\]

Spezialfall $i(0)=0$:
\[
\cbl{i(t) = \frac{U}{R}\left(1-\e^{-t/\tau}\right)}
\]

% -------------------------------------------------
\subsection{Beispiel 2: R--C-Kreis (Spannungsanstieg am Kondensator)}

KVL:
\[
\cbl{u_R + u_C = U,\quad u_R = Ri,\quad i = C\frac{du_C}{dt}}
\]

\[
RC\frac{du_C}{dt}+u_C = U
\qquad
\cbl{\tau = RC}
\]

Lösung:
\[
\cbl{u_C(t) = U + \big(u_C(0)-U\big)\e^{-t/\tau}}
\]

Spezialfall $u_C(0)=0$:
\[
\cbl{u_C(t) = U\left(1-\e^{-t/\tau}\right)}
\]

% -------------------------------------------------
\subsection{Integrierender Faktor (allgemeine 1. Ordnung)}

Für:
\[
\cbl{\frac{dy}{dt} + p(t) y = q(t)}
\]

Integrierender Faktor:
\[
\cbl{\mu(t) = \e^{\int p(t)\,dt}}
\]

Dann:
\[
\cbl{y(t) = \frac{1}{\mu(t)}\left(\int \mu(t) q(t)\,dt + C\right)}
\]

% -------------------------------------------------
\subsection{DGL 2. Ordnung im Zeitbereich}

\subsubsection{Standardform}

\[
\cbl{\frac{d^2y}{dt^2} + 2\zeta\omega_0 \frac{dy}{dt} + \omega_0^2 y = K\omega_0^2 u(t)}
\]

Parameter:
\[
\cbl{\omega_0 = \sqrt{\frac{1}{LC}}}
\qquad
\cbl{\zeta = \frac{R}{2}\sqrt{\frac{C}{L}}}
\]

\subsubsection{Homogene Lösung über charakteristische Gleichung}

Ansatz:
\[
y_h(t)=\e^{st}
\]

Charakteristische Gleichung:
\[
\cbl{s^2 + 2\zeta\omega_0 s + \omega_0^2 = 0}
\]

Pole:
\[
\cbl{s_{1,2} = -\zeta\omega_0 \pm \omega_0\sqrt{\zeta^2-1}}
\]

\subsubsection{Fälle}

Unterdämpft ($0<\zeta<1$):
\[
\cbl{y_h(t)=\e^{-\zeta\omega_0 t}\left(C_1\cos(\omega_d t)+C_2\sin(\omega_d t)\right)}
\]
\[
\cbl{\omega_d = \omega_0\sqrt{1-\zeta^2}}
\]

Kritisch ($\zeta=1$):
\[
\cbl{y_h(t)=(C_1+C_2 t)\e^{-\omega_0 t}}
\]

Überdämpft ($\zeta>1$):
\[
\cbl{y_h(t)=C_1\e^{s_1 t}+C_2\e^{s_2 t}}
\]

\subsubsection{Partikuläre Lösung}

Typische Prüfungsinputs:
\begin{itemize}
    \item Konstante $u(t)=U$
    \item Sprung $u(t)=U\cdot 1(t)$
    \item Sinus $u(t)=\hat U\sin(\omega t)$ (Frequenzantwort)
\end{itemize}

Für $u(t)=U$ ist meist $y_p$ konstant.

% -------------------------------------------------
\subsection{Beispiel 3: Serien-RLC (typische Form)}

Aus KVL:
\[
u(t)=Ri(t)+L\frac{di}{dt}+\frac{1}{C}\int i(t)\,dt
\]

Differenziert:
\[
\cbl{\frac{d^2 i}{dt^2} + \frac{R}{L}\frac{di}{dt} + \frac{1}{LC}i = \frac{1}{L}\frac{du}{dt}}
\]

Typischer Prüfungsfall: Spannungssprung $u(t)=U$ $\Rightarrow du/dt = U\delta(t)$.
Dann lösen über Laplace am saubersten.

% -------------------------------------------------
\subsection{Laplace-Methode (Standardworkflow)}

\subsubsection{Workflow}

\begin{enumerate}
    \item DGL aufstellen
    \item Laplace anwenden
    \item Anfangswerte einsetzen
    \item Nach $Y(s)$ auflösen
    \item Partialbruchzerlegung
    \item Rücktransformation
\end{enumerate}

\subsubsection{Merksätze}

\[
\cbl{\mathcal{L}\left\{\frac{dy}{dt}\right\}=sY(s)-y(0)}
\qquad
\cbl{\mathcal{L}\left\{\frac{d^2y}{dt^2}\right\}=s^2Y(s)-sy(0)-y'(0)}
\]

\subsubsection{Beispiel: 1. Ordnung per Laplace}

\[
L\frac{di}{dt}+Ri=U
\]

Laplace:
\[
LsI(s)-L i(0)+RI(s)=\frac{U}{s}
\]

\[
\cbl{I(s)=\frac{\frac{U}{s}+L i(0)}{Ls+R}}
\]

Rücktransformation ergibt:
\[
i(t)=\frac{U}{R}+\left(i(0)-\frac{U}{R}\right)\e^{-tR/L}
\]

% -------------------------------------------------
\subsection{Stückweise DGL bei Schaltvorgängen}

Typisch bei:
\begin{itemize}
    \item Gleichstromsteller (Ein/Aus-Phasen)
    \item Wechselrichter (Halbperioden)
    \item Wechselstromsteller (Leit-/Sperrintervalle)
\end{itemize}

\subsubsection{Vorgehen}

\begin{enumerate}
    \item Zeitintervalle definieren ($[t_0,t_1)$, $[t_1,t_2)$, \dots)
    \item Für jedes Intervall DGL aufstellen
    \item Allgemeine Lösung je Intervall berechnen
    \item Übergangsbedingungen anwenden:
    \[
    \cbl{i_L(t_1^-)=i_L(t_1^+)}
    \qquad
    \cbl{u_C(t_1^-)=u_C(t_1^+)}
    \]
    \item Periodenbedingung falls periodisch:
    \[
    \cbl{y(t)=y(t+T)}
    \]
\end{enumerate}

\crd{\textrightarrow\ Ohne Übergangs- und Periodenbedingungen ist die Lösung wertlos.}

% -------------------------------------------------
\subsection{Typische Prüfungs-Tricks}

\begin{itemize}
    \item Endwert sofort aus stationärem Zustand bestimmen ($dy/dt=0$).
    \item Zeitkonstante als $\tau=L/R$ oder $\tau=RC$ identifizieren.
    \item Stetigkeit von $i_L$ und $u_C$ konsequent nutzen.
    \item Einheitencheck: $\tau$ muss Sekunden sein.
\end{itemize}

\subsection{Typische Fehlerquellen}

\begin{itemize}
    \item Anfangswerte vergessen oder falsch eingesetzt.
    \item Exponentialfaktor falsches Vorzeichen.
    \item Stetigkeit verletzt ($i_L$ springt, $u_C$ springt).
    \item Stückweise Lösung ohne Randbedingungen.
\end{itemize}

\subsection{Minimal-Checkliste für Aufgaben}

\begin{enumerate}
    \item DGL-Typ identifizieren (1. Ordnung / 2. Ordnung / stückweise).
    \item $\tau$, $\omega_0$, $\zeta$ bestimmen.
    \item Homogen + Partikulär (oder Laplace) lösen.
    \item Konstanten mit Anfangs- und Übergangsbedingungen bestimmen.
    \item Grenzwerte $t\to 0$ und $t\to\infty$ prüfen.
\end{enumerate}
